\subsection{Parallelizing TKO (PTKO)}
\label{sec:ptko}

Sequential TKO starts with $minUtility = 1$ and raises it only as itemsets accumulate in the heap.
On large datasets this cold-start causes extensive exploration of low-utility branches before the
threshold becomes meaningful.
Our implementation \textbf{PTKO} addresses this with two complementary techniques.

\subsubsection{Technique 1: Speculative Initial Threshold}

After building all single-item utility lists (at the end of the second DB scan), the exact
$SumIutils$ of every 1-item utility list is known.
The $k$-th largest value among these is a valid lower bound on the final $minUtility$:
any itemset not reaching this bound cannot appear in the top-$k$ result.

\begin{algorithm}[H]
\caption{Speculative Threshold Bootstrap}
\begin{algorithmic}[1]
\Require Utility lists $\mathcal{L} = \{UL_1, \ldots, UL_m\}$, parameter $k$
\State $S \leftarrow \text{sort}(\{UL_i.SumIutils\},\;\text{descending})$
\If{$|\mathcal{L}| \geq k$}
    \State $minUtility \leftarrow S[k-1]$ \Comment{$k$-th largest single-item utility}
\EndIf
\end{algorithmic}
\end{algorithm}

This eliminates the cold-start overhead: workers begin with a threshold close to the final
answer and prune most low-utility branches from the very first iteration.
Note that only the single-item \emph{writeOut} check is affected; the upper-bound pruning
condition $SumIutils + SumRutils \geq minUtility$ still allows recursion into items whose
individual utility is below the threshold if their extension potential is sufficient.

\subsubsection{Technique 2: Top-Level Fork/Join with Atomic Threshold}

\paragraph{Correctness observation.}
In TKO's DFS, items are sorted ascending by TWU.
When the search processes top-level item $X_i$, the extended utility lists
$\widetilde{UL}(X_i, X_j)$ for all $j > i$ are \emph{freshly constructed} (new allocations)
and are private to the branch rooted at $X_i$.
No two top-level branches share mutable data --- only two primitives are shared:

\begin{enumerate}
  \item $minUtility$: the pruning threshold (raised monotonically, never lowered).
  \item $topKHeap$: the result collection (push/pop operations).
\end{enumerate}

\paragraph{Worker pool and depth cutoff.}
Top-level items are partitioned into $W$ contiguous ranges (one per worker goroutine).
Each goroutine explores its range sequentially, recursing into sub-trees without spawning
further goroutines.
This \emph{depth cutoff at 0} avoids goroutine explosion while still exposing enough
parallelism to utilize all cores.

\begin{algorithm}[H]
\caption{PTKO: Top-Level Fork/Join with Atomic Threshold}
\begin{algorithmic}[1]
\Require $\mathcal{L} = [X_0, \ldots, X_{n-1}]$ (sorted ascending by TWU), workers $W$
\State $chunk \leftarrow \lceil n / W \rceil$
\For{$w \leftarrow 0$ \textbf{to} $W-1$} \textbf{in parallel} \Comment{Fork}
    \For{$i \leftarrow w \cdot chunk$ \textbf{to} $\min((w{+}1)\cdot chunk,\,n)-1$}
        \State $minU \leftarrow \textsc{AtomicLoad}(minUtility)$
        \If{$X_i.SumIutils \geq minU$}
            \State \textsc{WriteOut}$(\emptyset,\; X_i)$ \Comment{mutex-protected heap push}
        \EndIf
        \If{$X_i.SumIutils + X_i.SumRutils \geq \textsc{AtomicLoad}(minUtility)$}
            \State $exULs \leftarrow [\widetilde{UL}(X_i, X_j) \mid j > i]$
            \State \textsc{Search}$([X_i.item],\; X_i,\; exULs)$ \Comment{sequential DFS}
        \EndIf
    \EndFor
\EndFor \Comment{Join — wait for all goroutines}
\end{algorithmic}
\end{algorithm}

\paragraph{Atomic threshold propagation.}
$minUtility$ is stored as an \texttt{atomic.Int64}.
\textsc{WriteOut} holds a mutex only long enough to push/trim the heap;
after releasing the mutex it raises $minUtility$ using a CAS retry loop:

\begin{algorithm}[H]
\caption{CAS Raise of Atomic $minUtility$}
\begin{algorithmic}[1]
\Procedure{RaiseMinUtil}{$newVal$}
\Repeat
    \State $old \leftarrow \textsc{AtomicLoad}(minUtility)$
    \If{$newVal \leq old$} \Return \EndIf
\Until{$\textsc{CAS}(minUtility,\; old,\; newVal)$}
\EndProcedure
\end{algorithmic}
\end{algorithm}

Because CAS is a \textbf{raise-only} operation, concurrent calls always converge to
the highest value seen so far.
Other goroutines observe the raised threshold on their next $\textsc{AtomicLoad}$ call
at the top of their pruning check --- no lock is needed for reads.

\paragraph{Challenge: unequal subtree sizes.}
Items are sorted by TWU ascending, so the last range (high-TWU items) tends to have
deeper, more expensive sub-trees.
Static partitioning may leave some workers idle while others are still busy.
A work-stealing task queue (channel-based) would improve balance at the cost of
implementation complexity; we leave this as future work.

\paragraph{Thread safety argument.}
\begin{itemize}
  \item \emph{Utility lists built in secondPass} are never mutated after construction
        (Elements slice is only appended to during secondPass, then frozen).
        Concurrent reads during \textsc{Construct} are safe.
  \item \emph{mapItemToTWU} is read-only after firstPass.
  \item \emph{kHeap} mutations are serialized by \texttt{sync.Mutex}.
  \item \emph{$minUtility$} is updated only via \textsc{RaiseMinUtil} (CAS, raise-only).
        Goroutines may briefly use a stale (lower) threshold, causing a few extra
        utility-list constructions that are harmless: their results pass through
        \textsc{WriteOut} and are evicted from the heap if below the final top-$k$ boundary.
\end{itemize}
